% GENERATED FILE. Edit canonical lessons, then run npm run generate:books.
% canonical-source-hash: fnv1a64:02d47384ee353b15
% canonical-lessons: AR-C238-mithl, AR-C238-tamaman, AR-C238-taban, AR-C238-taqriban, AR-C238-khususan

\chapter{Like, Completely, Of course, Almost, Especially}
\label{ch:ar-like-completely-of-course-almost-especially}
\hlchaptermodality{238}

\begin{chapteropening}
\textbf{By the end of this chapter:} \emph{I can say like, completely, of course, almost and especially.}

Complete the last lesson of chapter 238: I can say like, completely, of course, almost and especially.
\end{chapteropening}

\section[\texorpdfstring{\ar{مثل} (mithl)}{مثل (mithl)}]{\ar{مثل} (mithl) — like, such as}

\label{lesson:AR-C238-mithl}



\begin{quote}
Before the new one: say the Arabic for morals, then the Arabic for personality.
\end{quote}

\subsection*{You'll want to know: \ar{مثل}}
\textbf{\ar{مثل}} — \emph{mithl} — \textquotedblleft{}like, such as\textquotedblright{}.

\begin{cousinweb}
\textbf{\ar{مثل}} (\emph{mithl}) is built on the native Arabic root \emph{m-th-l}, \textquotedblleft{}to be like\textquotedblright{} (Lane). \emph{Mithl} is \textquotedblleft{}a like, a match\textquotedblright{}.
\end{cousinweb}

\begin{grammarlens}[title={What you've built}]
One more word for this chapter.
\end{grammarlens}

\subsection*{Your turn}
\begin{itemize}
\raggedright
  \item \emph{Say it:} \emph{mithl}
  \item \emph{Say it:} \emph{mithl}, once more
  \item \emph{Say it:} say \emph{mithl}
  \item \emph{Recall:} say the Arabic for morals, then the Arabic for personality, then say \emph{mithl} again
\end{itemize}

\begin{tcolorbox}[breakable,colback=teal!4,colframe=teal!35!black,title={Before you move on}]
What does \ar{مثل} mean? (\textquotedblleft{}Like, such as\textquotedblright{}.) Say it once more.
\end{tcolorbox}
\section[\texorpdfstring{\ar{تماما} (tamāman)}{تماما (tamāman)}]{\ar{تماما} (tamāman) — completely}

\label{lesson:AR-C238-tamaman}



\begin{quote}
Before the new one: say the Arabic for personality, then the Arabic for like.
\end{quote}

\subsection*{You'll want to know: \ar{تماما}}
\textbf{\ar{تماما}} — \emph{tamāman} — \textquotedblleft{}completely\textquotedblright{}.

\begin{cousinweb}
\textbf{\ar{تماما}} (\emph{tamāman}) is built on the native Arabic root \emph{t-m-m}, \textquotedblleft{}to be complete\textquotedblright{} (Lane). \emph{Tamāman} is \emph{tamām}, \textquotedblleft{}completeness\textquotedblright{}, with the adverb ending \emph{-an}.
\end{cousinweb}

\begin{grammarlens}[title={What you've built}]
One more word for this chapter.
\end{grammarlens}

\subsection*{Your turn}
\begin{itemize}
\raggedright
  \item \emph{Say it:} \emph{tamāman}
  \item \emph{Say it:} \emph{tamāman}, once more
  \item \emph{Say it:} say \emph{tamāman}
  \item \emph{Recall:} say the Arabic for personality, then the Arabic for like, then say \emph{tamāman} again
\end{itemize}

\begin{tcolorbox}[breakable,colback=teal!4,colframe=teal!35!black,title={Before you move on}]
What does \ar{تماما} mean? (\textquotedblleft{}Completely\textquotedblright{}.) Say it once more.
\end{tcolorbox}
\section[\texorpdfstring{\ar{طبعا} (ṭabʿan)}{طبعا (ṭabʿan)}]{\ar{طبعا} (ṭabʿan) — of course}

\label{lesson:AR-C238-taban}



\begin{quote}
Before the new one: say the Arabic for like, then the Arabic for completely.
\end{quote}

\subsection*{You'll want to know: \ar{طبعا}}
\textbf{\ar{طبعا}} — \emph{ṭabʿan} — \textquotedblleft{}of course\textquotedblright{}.

\begin{cousinweb}
\textbf{\ar{طبعا}} (\emph{ṭabʿan}) is built on the native Arabic root \emph{ṭ-b-ʿ}, \textquotedblleft{}to seal, to stamp\textquotedblright{} (Lane). \emph{Ṭabʿan} is \emph{ṭabʿ}, \textquotedblleft{}nature\textquotedblright{}, with the adverb ending \emph{-an}: \textquotedblleft{}by nature, of course\textquotedblright{}.
\end{cousinweb}

\begin{grammarlens}[title={What you've built}]
One more word for this chapter.
\end{grammarlens}

\subsection*{Your turn}
\begin{itemize}
\raggedright
  \item \emph{Say it:} \emph{ṭabʿan}
  \item \emph{Say it:} \emph{ṭabʿan}, once more
  \item \emph{Say it:} say \emph{ṭabʿan}
  \item \emph{Recall:} say the Arabic for like, then the Arabic for completely, then say \emph{ṭabʿan} again
\end{itemize}

\begin{tcolorbox}[breakable,colback=teal!4,colframe=teal!35!black,title={Before you move on}]
What does \ar{طبعا} mean? (\textquotedblleft{}Of course\textquotedblright{}.) Say it once more.
\end{tcolorbox}
\section[\texorpdfstring{\ar{تقريبا} (taqrīban)}{تقريبا (taqrīban)}]{\ar{تقريبا} (taqrīban) — almost, approximately}

\label{lesson:AR-C238-taqriban}



\begin{quote}
Before the new one: say the Arabic for completely, then the Arabic for of course.
\end{quote}

\subsection*{You'll want to know: \ar{تقريبا}}
\textbf{\ar{تقريبا}} — \emph{taqrīban} — \textquotedblleft{}almost, approximately\textquotedblright{}.

\begin{cousinweb}
\textbf{\ar{تقريبا}} (\emph{taqrīban}) is built on the native Arabic root \emph{q-r-b}, \textquotedblleft{}to be near\textquotedblright{} (Lane). \emph{Taqrīban} is the verbal noun of Form II, \textquotedblleft{}bringing near\textquotedblright{}, with the adverb ending \emph{-an}.
\end{cousinweb}

\begin{grammarlens}[title={What you've built}]
One more word for this chapter.
\end{grammarlens}

\subsection*{Your turn}
\begin{itemize}
\raggedright
  \item \emph{Say it:} \emph{taqrīban}
  \item \emph{Say it:} \emph{taqrīban}, once more
  \item \emph{Say it:} say \emph{taqrīban}
  \item \emph{Recall:} say the Arabic for completely, then the Arabic for of course, then say \emph{taqrīban} again
\end{itemize}

\begin{tcolorbox}[breakable,colback=teal!4,colframe=teal!35!black,title={Before you move on}]
What does \ar{تقريبا} mean? (\textquotedblleft{}Almost, approximately\textquotedblright{}.) Say it once more.
\end{tcolorbox}
\section[\texorpdfstring{\ar{خصوصا} (khuṣūṣan)}{خصوصا (khuṣūṣan)}]{\ar{خصوصا} (khuṣūṣan) — especially}

\label{lesson:AR-C238-khususan}



\begin{quote}
Before the new one: say the Arabic for of course, then the Arabic for almost.
\end{quote}

\subsection*{You'll want to know: \ar{خصوصا}}
\textbf{\ar{خصوصا}} — \emph{khuṣūṣan} — \textquotedblleft{}especially\textquotedblright{}.

\begin{cousinweb}
\textbf{\ar{خصوصا}} (\emph{khuṣūṣan}) is built on the native Arabic root \emph{kh-ṣ-ṣ}, \textquotedblleft{}to single out\textquotedblright{} (Lane). \emph{Khuṣūṣan} is \emph{khuṣūṣ}, \textquotedblleft{}singling out\textquotedblright{}, with the adverb ending \emph{-an}.
\end{cousinweb}

\begin{grammarlens}[title={What you've built}]
One more word for this chapter.
\end{grammarlens}

\subsection*{Your turn}
\begin{itemize}
\raggedright
  \item \emph{Say it:} \emph{khuṣūṣan}
  \item \emph{Say it:} \emph{khuṣūṣan}, once more
  \item \emph{Say it:} say \emph{khuṣūṣan}
  \item \emph{Recall:} say the Arabic for of course, then the Arabic for almost, then say \emph{khuṣūṣan} again
\end{itemize}

\begin{tcolorbox}[breakable,colback=teal!4,colframe=teal!35!black,title={Before you move on}]
What does \ar{خصوصا} mean? (\textquotedblleft{}Especially\textquotedblright{}.) Say it once more.
\end{tcolorbox}
